\documentclass[11pt]{amsart}

\usepackage[T1]{fontenc}
\usepackage{lmodern}
\usepackage{microtype}
\usepackage{amsmath,amssymb}
\usepackage{booktabs}
\usepackage[hidelinks]{hyperref}
\hypersetup{
  pdftitle={Nine additional counterexamples to Conjecture 5 of Sondow, Nicholson, and Noe}
}

\newtheorem{theorem}{Theorem}

\title[Additional Ramanujan-prime exceptions]
{Nine additional counterexamples to Conjecture~5 of Sondow, Nicholson, and Noe}

\date{}

\subjclass[2020]{Primary 11N05; Secondary 11A41}
\keywords{Ramanujan prime, prime-counting function, explicit estimate,
computational number theory}

\begin{document}

\begin{abstract}
Let $R_n$ be the $n$th Ramanujan prime and put
$\rho(n)=\pi(R_n)$.  The $m\geq20$ clause of Conjecture~5 of
Sondow, Nicholson, and Noe asserts that
$\rho(mn)\leq m\rho(n)$ for every $n\geq2$.
Axler later exhibited the exception $(m,n)=(38,9)$ and asserted that it
was unique.  We give nine additional counterexamples and prove that, for
$m\geq20$, the complete exception set is
\[
 \{(m,9):m\in\{22,23,24,25,26,32,33,36,37,38\}\}.
\]
Consequently, the least valid threshold is $n=10$ for exactly these ten
values of $m$, and is $n=2$ for every other $m\geq20$.
\end{abstract}

\maketitle

All logarithms are natural.  The $n$th Ramanujan prime is the least
positive integer $R_n$ such that
\[
 \pi(x)-\pi(x/2)\geq n \qquad (x\geq R_n),
\]
and $p_k$ denotes the $k$th prime.  Sondow, Nicholson, and Noe
\cite[Conjecture~5]{SNN} conjectured, in particular, that
\begin{equation}\label{eq:conjecture}
 \rho(mn)\leq m\rho(n) \qquad (m\geq20,\ n\geq2).
\end{equation}
Axler subsequently asserted that the only exception to the full
conjecture is $(38,9)$ \cite[Theorem~1.7]{Axler17}; the paper later
appeared as \cite{Axler20}.

\begin{theorem}\label{thm:main}
For integers $m\geq20$ and $n\geq2$,
\[
 \rho(mn)>m\rho(n)
 \quad\Longleftrightarrow\quad
 n=9\ \text{ and }\
 m\in E,
\]
where
\[
 E=\{22,23,24,25,26,32,33,36,37,38\}.
\]
Equivalently, if
\[
 N_{\min}(m)=\min\{N\geq2:\rho(mn)\leq m\rho(n)
                         \text{ for every }n\geq N\},
\]
then
\[
 N_{\min}(m)=
 \begin{cases}
  10,&m\in E,\\
  2,&m\geq20\text{ and }m\notin E.
 \end{cases}
\]
\end{theorem}

\begin{proof}
We first verify the infinite tail.  Put
\[
 a=1.472,\qquad b=2.51,\qquad c=\log2,
 \qquad \phi(n)=a\log\log n+b.
\]
Axler's estimates \cite[(5.5) and (5.9)]{Axler17} reduce the desired
inequality, for $m\geq20$ and $n\geq1245$, to the positivity of a
quantity $W_m(n)$ satisfying
\begin{align*}
W_m(n)>{}&\log m\left(c+\frac{c}{\log n}
 -\frac{\phi(n)}{\log^2n}-\frac{\phi(n)}{\log^3n}\right)
 -\frac{\phi(n)}{\log n}\\
&-\frac{\phi(n)(\log\log n-c)}{\log^2n}\\
&+\frac{(\log m+\tfrac12)c-0.565}{\log(mn)}\\
&+\frac{(c\log m-0.565)\log\log n}{\log^2(mn)}.
\end{align*}
The last two terms are positive for $m\geq20$.  Define
\[
 q(n)=0.03+\frac{c}{\log n}
      -\frac{\phi(n)}{\log^2n}-\frac{\phi(n)}{\log^3n}.
\]
Let $t_0=\log1245$, $\Phi(t)=a\log t+b$, and
\[
 Q(t)=0.03t^3+ct^2-\Phi(t)(t+1).
\]
Then $Q(t)=t^3q(e^t)$ and
\[
 Q''(t)=0.18t+2c-\frac{a}{t}+\frac{a}{t^2}>2.46
 \qquad(t\geq t_0).
\]
Moreover, $Q'(t_0)>7.36$ and $Q(t_0)>2.15$, so $q(n)>0$ for
$n\geq1245$.  These numerical inequalities follow, for example, from
\[
 7.126<t_0<7.127,\qquad
 0.6931<c<0.6932,\qquad
 1.963<\log t_0<1.964.
\]

The functions
\[
 u_1(t)=\frac{\Phi(t)}{t},
 \qquad
 u_2(t)=\frac{\Phi(t)(\log t-c)}{t^2}
\]
are decreasing for $t\geq t_0$.  Indeed,
$u_1'(t)=(a-\Phi(t))/t^2<0$, while
\[
 t^3u_2'(t)=H(t):=a(\log t-c)+\Phi(t)
                    -2\Phi(t)(\log t-c),
\]
with
\[
 H'(t)=\frac{2(a-a(\log t-c)-\Phi(t))}{t}<0
 \quad\text{and}\quad H(t_0)<-6.44.
\]
It follows that
\begin{align*}
 W_m(n)
 &>(c-0.03)\log20-u_1(t_0)-u_2(t_0)\\
 &>1.09.
\end{align*}
Thus
\begin{equation}\label{eq:tail}
 \rho(mn)\leq m\rho(n)
 \qquad(m\geq20,\ n\geq1245).
\end{equation}

It remains to consider $2\leq n\leq1244$.  Define
\[
 \gamma(x)=
 \frac{c+c/\log x+0.565/\log^2x}
      {\log x+\log\log x-c-c/\log x}.
\]
Axler's upper estimate \cite[Theorem~1.2 and (3.4)--(3.5)]{Axler17}
gives
\begin{equation}\label{eq:upper}
 \rho(k)<2k(1+\gamma(k)) \qquad(k\geq5225).
\end{equation}
The numerator of $\gamma$ is positive and decreasing, and its denominator
is positive and increasing; hence $\gamma$ is decreasing on this range.
An exact computation gives
\begin{equation}\label{eq:epsilon}
 \min_{2\leq n\leq1244}
 \left(\frac{\rho(n)}{2n}-1\right)
 =\frac{101}{1244},
\end{equation}
attained uniquely at $n=1244$, where
$R_{1244}=24169$ and $\rho(1244)=2690$.  Also,
\[
 \gamma(5225)
 <\frac{0.694+0.694/8.56+0.565/8.56^2}
         {8.56+2.14-0.694-0.694/8.56}
 <0.079<\frac{101}{1244}.
\]
Therefore, if $2\leq n\leq1244$ and $mn\geq5225$, then
\[
 \rho(mn)
 <2mn(1+\gamma(mn))
 <m\rho(n).
\]

Only
\[
 \mathcal F=\{(m,n):m\geq20,\ 2\leq n\leq1244,\ mn<5225\}
\]
remains.  Since necessarily $n\leq261$,
\[
 \#\mathcal F
 =\sum_{n=2}^{261}
   \left(\left\lfloor\frac{5224}{n}\right\rfloor-19\right)
 =21803.
\]
For an integer $x\geq0$, put
\[
 d(x)=\pi(x)-\pi(\lfloor x/2\rfloor).
\]
By \cite[Theorem~3]{SNN}, $R_j<p_{3j}$ for every $j$.  Hence a sieve
through
\[
 B=p_{15672}=172069
\]
contains every $R_j$ with $j\leq5224$.  Since the function
$\pi(x)-\pi(x/2)$ changes only at integer points,
\begin{equation}\label{eq:lastbad}
 R_j=1+\max\{0\leq x<B:d(x)<j\}.
\end{equation}
This is the last-deficit computation described in
\cite[Appendix]{SNN}.

Applying \eqref{eq:lastbad} and checking all pairs in $\mathcal F$ in
integer arithmetic gives exactly the following failures.  Here
$R_9=71$ and $\rho(9)=20$.
\begin{center}
\renewcommand{\arraystretch}{1.05}
\begin{tabular}{@{}rrrrrr@{}}
\toprule
$m$ & $9m$ & $R_{9m}$ & $\rho(9m)$ & $20m$ & excess \\
\midrule
22 & 198 & 3167 & 448 & 440 & 8 \\
23 & 207 & 3319 & 467 & 460 & 7 \\
24 & 216 & 3467 & 486 & 480 & 6 \\
25 & 225 & 3607 & 504 & 500 & 4 \\
26 & 234 & 3761 & 523 & 520 & 3 \\
32 & 288 & 4789 & 644 & 640 & 4 \\
33 & 297 & 4967 & 664 & 660 & 4 \\
36 & 324 & 5477 & 723 & 720 & 3 \\
37 & 333 & 5651 & 743 & 740 & 3 \\
38 & 342 & 5807 & 762 & 760 & 2 \\
\bottomrule
\end{tabular}
\end{center}
The last row is Axler's previously known exception; the first nine are
additional.  There are no other failures in $\mathcal F$.  Together with
\eqref{eq:tail} and \eqref{eq:upper}, this proves the classification.
The formula for $N_{\min}(m)$ follows immediately.
\end{proof}

\begin{thebibliography}{9}

\bibitem{Axler17}
C.~Axler,
\emph{On the number of primes up to the $n$th Ramanujan prime},
arXiv:1711.04588v1 [math.NT], 2017.

\bibitem{Axler20}
C.~Axler,
\emph{On Ramanujan primes},
Funct. Approx. Comment. Math. \textbf{63} (2020), no.~1, 67--93,
\href{https://doi.org/10.7169/facm/1824}{doi:10.7169/facm/1824}.

\bibitem{SNN}
J.~Sondow, J.~W.~Nicholson, and T.~D.~Noe,
\emph{Ramanujan primes: bounds, runs, twins, and gaps},
J. Integer Seq. \textbf{14} (2011), no.~6, Article 11.6.2, 11~pp.

\end{thebibliography}

\end{document}